\input{uni}

\begin{document}

\setstretch{1.5}
\noindent
\begin{minipage}[t]{0.01\textwidth}
\includegraphics{/home/stanik/programmer/++Programmer/5_2/LSTU.jpg}
\end{minipage}%
\begin{minipage}[c]{1.1\linewidth}
    \centering
    {\small
    МИНИСТЕРСТВО НАУКИ И ВЫСШЕГО ОБРАЗОВАНИЯ РФ \\[1em]
    ФЕДЕРАЛЬНОЕ ГОСУДАРСТВЕННОЕ БЮДЖЕТНОЕ \\
    ОБРАЗОВАТЕЛЬНОЕ УЧРЕЖДЕНИЕ ВЫСШЕГО ОБРАЗОВАНИЯ \\
    «ЛИПЕЦКИЙ ГОСУДАРСТВЕННЫЙ ТЕХНИЧЕСКИЙ УНИВЕРСИТЕТ»
    }
\end{minipage}%

\vspace{3em}

\noindent
\hspace{1em}%
\begin{minipage}[t]{0.8\textwidth}
\textbf{Институт} \hspace{2em} компьютерных наук \\
\textbf{Кафедра} \hspace{2em} автоматизированных систем управления
\end{minipage}


\vspace{6em}

\begin{center}
\textbf{Индивидуальное домашнее задание} \\
по дисциплине \\ <<Управление ИТ-проектами>>
\end{center}

\vspace{6em}

\noindent
\begin{tabular*}{\textwidth}{@{\extracolsep{\fill}} l c r }
    Студент М-РИТ-25-1 & \makebox[4cm]{\hrulefill} & Станиславчук С.М. \\[2em]
Руководитель & \\
К.т.н., доцент & \makebox[4cm]{\hrulefill} & Ведищев В.В. \\
\end{tabular*}

\vfill
\begin{center}
    Липецк, 2026
\end{center}

\thispagestyle{empty}

\newpage
\tableofcontents
\newpage

\section{Разработка прототипа ИТ-проекта}

\subsection{Используемые средства для разработки прототипа}

Разработка прототипа программного комплекса <<CorpCom>> выполнена с~использованием языка Python~3.13 и~стека библиотек, обеспечивающих сетевое взаимодействие в~реальном времени, графический интерфейс и~потоковую обработку звука. Выбор инструментов обоснован требованиями к~прототипу: работа в~локальной сети предприятия без~внешних облачных сервисов, минимальные задержки голосовой связи, кроссплатформенность клиентской части.

\subsubsection{Язык программирования}

В~качестве основного языка разработки выбран Python (версия~3.13). Решение мотивировано следующими факторами:

\begin{itemize}
  \item \textbf{Скорость прототипирования.} Python позволяет реализовать WebSocket-сервер, клиентский GUI и~голосовой движок в~считанные часы благодаря лаконичному синтаксису и~богатой экосистеме библиотек.
  \item \textbf{Асинхронная модель.} Встроенный модуль \texttt{asyncio} обеспечивает конкурентную обработку десятков WebSocket-соединений на~одном потоке без~накладных расходов на~многопоточную синхронизацию.
  \item \textbf{Кроссплатформенность.} Код выполняется на~любой ОС с~установленным интерпретатором Python (Linux, macOS, Windows), что~упрощает развёртывание прототипа в~разнородной корпоративной среде.
\end{itemize}

\subsubsection{Сетевое взаимодействие}

Для реализации транспортного уровня протокола выбрана библиотека \texttt{websockets} (версия~$\ge$~12.0). Библиотека предоставляет:

\begin{itemize}
  \item асинхронный сервер (\texttt{websockets.serve}) с~встроенной поддержкой \texttt{ping/pong} для~обнаружения разрывов соединения;
  \item асинхронный клиент (\texttt{websockets.connect}) с~контекстным менеджером, автоматически закрывающим соединение;
  \item потоковый приём сообщений через \texttt{async for} по~сокету;
  \item настраиваемый лимит размера входящих сообщений (\texttt{max\_size}) — критично для~передачи изображений.
\end{itemize}

Протокол обмена --- JSON поверх WebSocket (RFC~6455). Все сообщения кодируются в~UTF-8 и~передаются как текстовые фреймы. Бинарные данные (изображения, PCM-аудио) кодируются в~Base64 и~вкладываются в~JSON-поле \texttt{data}. Такой подход упрощает отладку: любое сообщение можно прочитать в~консоли браузера или~Wireshark без~десериализации бинарного потока.

\subsubsection{Графический интерфейс}

Клиентский GUI реализован на~библиотеке \texttt{tkinter} (стандартный модуль Python, обвязка над~Tcl/Tk). Выбор обоснован:

\begin{itemize}
  \item отсутствием внешних зависимостей --- tkinter входит в~стандартную поставку Python;
  \item достаточной функциональностью для~прототипа: текстовые поля, кнопки, скроллируемый чат, диалоги выбора файлов, превью изображений;
  \item событийной моделью, совместимой с~асинхронным сетевым кодом через очередь \texttt{queue.Queue} и~периодический опрос \texttt{root.after()}.
\end{itemize}

GUI разделён на~смысловые панели: <<Подключение>> (адрес, порт, имя, комната, токен), <<Голос в~реальном времени>> (микрофон, PTT, приём), <<Чат>> (история сообщений, поле ввода) и~<<Медиа>> (отправка изображения на~проектор).

Для клиента проектора также используется tkinter с~полноэкранным режимом (\texttt{attributes("-fullscreen")}) и~автоматическим масштабированием изображения под~размер канваса.

\subsubsection{Голосовой движок}

Потоковая обработка звука реализована на~библиотеках \texttt{sounddevice} (захват и~воспроизведение PCM через PortAudio) и~\texttt{numpy} (работа с~массивами семплов). Параметры аудио:

\begin{itemize}
  \item частота дискретизации: 16~000 Гц;
  \item разрядность: 16~бит;
  \item канальность: моно;
  \item размер блока: 320 семплов (20~мс).
\end{itemize}

Выбор 16~кГц моно обусловлен компромиссом между качеством речи и~пропускной способностью сети: каждый 20-миллисекундный блок занимает 640~байт (без~учёта Base64-кодирования и~JSON-обёртки), что~даёт поток около~32~КБ/с — приемлемо даже для~Wi-Fi сетей предприятия.

Захват звука реализован через \texttt{sd.InputStream} с~callback-функцией, которая передаёт блоки байт в~асинхронную очередь \texttt{asyncio.Queue}. Воспроизведение --- через \texttt{sd.OutputStream} с~непрерывным callback-чтением из~потокобезопасной очереди \texttt{queue.Queue}. При~переполнении очереди старые блоки отбрасываются --- это предотвращает накопление задержки.

\subsubsection{Обработка изображений}

Для работы с~изображениями используется библиотека \texttt{Pillow} (PIL, версия~$\ge$~10.0). В~клиенте Pillow применяется для:

\begin{itemize}
  \item чтения файлов изображений с~диска в~бинарный буфер (\texttt{Image.open});
  \item создания уменьшенного превью (320$\times$180) для~отправителя;
  \item масштабирования полученного изображения под~размер канваса проектора (\texttt{Image.thumbnail} с~фильтром \texttt{LANCZOS}).
\end{itemize}

\subsubsection{Виртуальное окружение и управление зависимостями}

Зависимости проекта зафиксированы в~файле \texttt{requirements.txt} и~устанавливаются в~изолированное виртуальное окружение \texttt{.venv/}. Скрипт \texttt{setup.sh} автоматизирует создание окружения и~установку пакетов, а~скрипт \texttt{corpcom\_exec.sh} обеспечивает запуск любого модуля с~корректным \texttt{PYTHONPATH} и~\texttt{PATH}, обходя проблему подмены \texttt{sys.executable} в~IDE Cursor.

\subsubsection{Контроль версий}

Исходный код размещён в~системе контроля версий Git, что~позволяет отслеживать изменения на~всех этапах разработки прототипа и~возвращаться к~предыдущим состояниям при~необходимости.

\subsection{Функционал прототипа}

Прототип ПК~<<CorpCom>> реализует три~ключевые функции, обеспечивающие полноценный обмен информацией внутри рабочей группы: текстовые сообщения, голосовой канал и~передачу файлов (изображений) на~общий экран.

\subsubsection{Текстовые сообщения}

Обмен текстовыми сообщениями реализован по~модели <<комната>> (chat~room). Пользователь подключается к~именованной комнате и~получает сообщения от~всех остальных участников этой~комнаты в~реальном времени.

\textbf{Протокол.} Текстовое сообщение передаётся как JSON-объект типа \texttt{MSG\_TEXT} с~полями: \texttt{room} (комната), \texttt{sender} (отправитель), \texttt{content} (текст). Сервер при~получении текстового сообщения выполняет:

\begin{enumerate}
  \item обрезку содержимого до~8000 символов (\texttt{content[:8000]});
  \item проверку на~пустое сообщение (пустые игнорируются);
  \item широковещательную рассылку всем пирам комнаты с~ролью \texttt{user}, исключая отправителя.
\end{enumerate}

\textbf{Интерфейс.} В~клиенте история чата отображается в~виджете \texttt{ScrolledText}, новые сообщения автоматически прокручиваются в~конец. Сообщения от~других участников маркируются префиксом \texttt{[имя]}, собственные --- префиксом \texttt{[я]}. Отправка производится по~клавише \texttt{Enter} или~кнопке <<Отправить>>.

\textbf{Системные уведомления.} Сервер автоматически генерирует системные сообщения (\texttt{MSG\_SYS}) при~подключении и~отключении участников, а~также рассылает актуальный список участников комнаты (\texttt{MSG\_CLIENT\_LIST}) после каждого изменения состава.

\subsubsection{Голосовой канал}

Голосовая связь реализована как потоковая передача PCM-аудио поверх WebSocket в~режиме реального времени. В~отличие от~традиционной VoIP-телефонии, голосовой канал CorpCom не~требует установки выделенного медиа-соединения (SIP/RTP) --- аудиоданные передаются через тот~же WebSocket-канал, что~и~текстовые сообщения, что~упрощает архитектуру и~обходит проблемы с~NAT/брандмауэрами.

\textbf{Режимы передачи.} Прототип поддерживает два~режима:

\begin{itemize}
  \item \textbf{Открытый микрофон} --- непрерывная передача голоса. Пользователь нажимает <<Включить микрофон>>, после~чего все захваченные PCM-блоки автоматически отправляются на~сервер и~ретранслируются остальным участникам комнаты. При~нажатии <<Выключить микрофон>> отправляется сигнал \texttt{MSG\_VOICE\_END}.
  \item \textbf{PTT (Push-to-Talk)} --- передача только при~удержании кнопки. При~нажатии кнопки PTT отправляется \texttt{MSG\_VOICE\_START} и~начинается захват; при~отпускании --- \texttt{MSG\_VOICE\_END} и~захват останавливается. Режим удобен в~условиях фонового шума и~при~участии в~совещании нескольких человек из~одного помещения.
\end{itemize}

Режимы взаимоисключающие: при~включённом открытом микрофоне кнопка PTT игнорируется.

\textbf{Формат аудио.} Каждый голосовой блок --- это 320 семплов PCM (16~бит, моно, 16~кГц), что~соответствует 20~мс звука. Блок кодируется в~Base64 и~отправляется как JSON-сообщение типа \texttt{MSG\_VOICE} с~полями \texttt{room}, \texttt{sender}, \texttt{seq} (порядковый номер) и~\texttt{data} (Base64-строка). Порядковый номер позволяет детектировать потерю блоков, хотя в~текущей версии прототипа восстановление пропущенных блоков не~реализовано.

\textbf{Приём.} На~стороне получателя голосовые блоки декодируются из~Base64, преобразуются в~массив \texttt{numpy.int16} и~подаются в~кольцевой буфер воспроизведения. \texttt{sd.OutputStream} в~режиме непрерывного callback забирает блоки из~буфера и~выводит на~аудиоустройство. При~переполнении буфера старые блоки отбрасываются, что~предотвращает неограниченный рост задержки.

Управление приёмом осуществляется галочкой <<Слышать собеседников>>. При~снятии галочки аудиопоток останавливается (вызов \texttt{OutputStream.stop()}), но~текстовые уведомления о~начале/конце голосовой передачи продолжают отображаться в~чате.

\subsubsection{Передача файлов (изображений на проектор)}

Третья функция прототипа --- передача графических файлов (изображений) на~клиент проектора для~демонстрации на~общем экране.

\textbf{Сценарий.} Пользователь нажимает кнопку <<Изображение на~проектор>>, выбирает файл через~стандартный диалог ОС (поддерживаются форматы PNG, JPEG, BMP, GIF, WebP), после~чего:

\begin{enumerate}
  \item Клиент проверяет размер файла (лимит 12~МБ). При~превышении выводится ошибка <<Файл слишком большой>>.
  \item Файл читается в~бинарный буфер, кодируется в~Base64 и~отправляется как JSON-сообщение типа \texttt{MSG\_IMAGE} с~полями \texttt{room}, \texttt{sender}, \texttt{filename}, \texttt{data}.
  \item Сервер рассылает сообщение всем пирам комнаты (роли \texttt{user} и~\texttt{projector}). Таким образом, отправитель и~другие пользователи также получают изображение --- это служит подтверждением отправки.
  \item Клиенты-проекторы при~получении \texttt{MSG\_IMAGE} декодируют Base64, открывают изображение через Pillow, масштабируют под~размер канваса и~отображают.
  \item Клиент-отправитель создаёт уменьшенное превью (320$\times$180) в~нижней части окна и~запись в~чате.
\end{enumerate}

\textbf{Полноэкранный режим.} Клиент проектора поддерживает переключение в~полноэкранный режим по~клавише \texttt{F11} и~выход из~него по~\texttt{Escape}. В~полноэкранном режиме скрываются панели управления и~строка состояния, изображение занимает всё пространство экрана.

\subsection{Запуск прототипа}

Запуск прототипа ПК~<<CorpCom>> выполняется в~три этапа: запуск сервера, запуск клиентов и~подключение к~нужной комнате. Все компоненты запускаются из~корневой директории проекта после однократного выполнения скрипта настройки \texttt{setup.sh}, создающего виртуальное окружение Python с~необходимыми зависимостями.

\subsubsection{Шаг 1. Сервер}

Сервер (\texttt{corpcom.server}) запускается первым, так~как клиенты и~проектор подключаются к~нему. На~хосте, выделенном под~сервер, выполняется:

\begin{verbatim}
  $ cd <корень проекта>
  $ ./run_server.sh
\end{verbatim}

Сервер стартует на~адресе \texttt{0.0.0.0} (все~сетевые интерфейсы) и~порту~8765 (значения по~умолчанию). При~успешном запуске в~консоль выводится:

\begin{verbatim}
  CorpCom server 0.0.0.0:8765 (token из CORPCOM_TOKEN)
\end{verbatim}

Сервер работает в~бесконечном цикле событий asyncio, ожидая входящие WebSocket-соединения. Для~завершения используется \texttt{Ctrl+C}.

При~необходимости изменить порт или~токен аутентификации, они задаются через переменные окружения:

\begin{verbatim}
  $ CORPCOM_PORT=9000 CORPCOM_TOKEN=my-safe-token ./run_server.sh
\end{verbatim}

Аналогично можно указать конкретный IP-адрес для~привязки (если хост имеет несколько интерфейсов):

\begin{verbatim}
  $ CORPCOM_HOST=192.168.1.10 ./run_server.sh
\end{verbatim}

\subsubsection{Шаг 2. Клиент}

На~компьютере каждого пользователя запускается клиентское приложение с~графическим интерфейсом:

\begin{verbatim}
  $ cd <корень проекта>
  $ ./run_client.sh
\end{verbatim}

Открывается окно клиента CorpCom. В~панели <<Подключение (LAN)>> заполняются поля:

\begin{itemize}
  \item \textbf{Сервер} --- IP-адрес хоста с~запущенным сервером (по~умолчанию \texttt{127.0.0.1} для~локального тестирования);
  \item \textbf{Порт} --- TCP-порт сервера (по~умолчанию~8765);
  \item \textbf{Имя} --- отображаемое имя пользователя в~чате (до~64 символов);
  \item \textbf{Комната} --- идентификатор комнаты. Все~участники, указавшие одинаковую комнату, будут обмениваться сообщениями и~голосом. По~умолчанию \texttt{general};
  \item \textbf{Токен} --- секретный токен аутентификации. Должен совпадать с~токеном, заданным при~запуске сервера (переменная \texttt{CORPCOM\_TOKEN}).
\end{itemize}

Нажатие кнопки \textbf{<<Подключить>>} инициирует WebSocket-соединение с~сервером. Клиент отправляет AUTH-сообщение с~токеном, ролью \texttt{user}, именем и~комнатой. При~успешной аутентификации:

\begin{itemize}
  \item статусная строка меняется на~<<В~сети: \textit{комната}>> с~указанием списка участников;
  \item в~чате появляется системное сообщение <<Подключено к~CorpCom>>;
  \item кнопка подключения меняется на~<<Отключить>>.
\end{itemize}

При~неверном токене сервер возвращает ошибку аутентификации, и~клиент показывает диалоговое окно с~текстом ошибки.

\subsubsection{Шаг 3. Клиент проектора (опционально)}

Если требуется демонстрация изображений на~общем экране, на~компьютере, подключённом к~проектору или~большому монитору, запускается:

\begin{verbatim}
  $ cd <корень проекта>
  $ ./run_projector.sh
\end{verbatim}

В~отличие от~обычного клиента, проектор подключается с~ролью \texttt{projector}. Поля <<Сервер>>, <<Порт>>, <<Комната>> и~<<Токен>> заполняются аналогично клиенту. После~нажатия кнопки <<Старт>> проектор переходит в~режим ожидания изображений. При~получении изображения оно автоматически отображается на~всём канвасе. Клавиша \texttt{F11} переводит окно в~полноэкранный режим.

\subsubsection{Шаг 4. Подключение к~нужной комнате}

Механизм комнат реализован на~стороне сервера: каждый пир при~аутентификации указывает имя комнаты, и~сервер маршрутизирует сообщения только~между пирами одной комнаты. Таким образом, несколько рабочих групп могут одновременно использовать один сервер, не~пересекаясь по~данным.

Для~подключения к~нужной комнате достаточно указать её~идентификатор в~поле <<Комната>> клиента перед~нажатием <<Подключить>>. Идентификатор комнаты --- произвольная строка до~64 символов. Рекомендуется использовать осмысленные имена, отражающие назначение комнаты: \texttt{general}, \texttt{lab}, \texttt{project-alpha}, \texttt{soveshchanie}.

\textbf{Пример сценария.} Три~участника проектной группы хотят провести совещание с~демонстрацией схем:

\begin{enumerate}
  \item Администратор запускает сервер: \texttt{CORPCOM\_TOKEN=meeting2026 ./run\_server.sh}
  \item Докладчик запускает клиент, указывает сервер, комнату \texttt{soveshchanie}, токен \texttt{meeting2026}, нажимает <<Подключить>>.
  \item Два~других участника запускают клиенты с~теми~же параметрами.
  \item На~компьютере проектора запускается \texttt{./run\_projector.sh} с~комнатой \texttt{soveshchanie}.
  \item Все~видят друг друга в~списке участников. Докладчик включает микрофон и~начинает отправлять изображения слайдов на~проектор. Остальные участники слушают доклад и~задают вопросы текстом в~чате.
\end{enumerate}

После~завершения совещания каждый участник нажимает <<Отключить>>, проектор --- <<Стоп>>, администратор завершает сервер по~\texttt{Ctrl+C}.

\end{document}
